\documentclass[12pt,a4paper]{article}

\usepackage[utf8]{inputenc}
\usepackage[T1]{fontenc}
\usepackage{lmodern}
\usepackage[french]{babel}
\usepackage{graphicx}
\usepackage{geometry}
\usepackage{fancyhdr}
\usepackage{tikz}
\usepackage{tabularx}
\usepackage{colortbl}
\usepackage{xcolor}
\usepackage[hidelinks]{hyperref}
\usepackage{longtable}
\usepackage{array}
\usepackage{enumitem}

\geometry{
    a4paper,
    top=3cm,
    bottom=3cm,
    left=2.5cm,
    right=2.5cm,
}

\pagestyle{empty}

\begin{document}

\pagestyle{plain}

\begin{titlepage}
    \centering

    \includegraphics[width=0.3\textwidth]{univ_rouen.jpg}

    \vspace{2cm}
    
    {\Large Master 1 GIL --- Technologies XML \\ \textcolor{blue!70!black}{\textbf{Projet : Service REST SEPA (ISO 20022) et Déploiement Conteneurisé}}}

    \vspace{2cm}
    
    \begin{center}
        \rule{\textwidth}{1.5pt}
        \vskip 0.8cm
        {\huge\bfseries RAPPORT TECHNIQUE DE PROJET}
        \vskip 0.8cm
        \rule{\textwidth}{1.5pt}
    \end{center}
    
    \vspace{2cm}
    
    \fbox{%
        \begin{minipage}{12cm}
            \textbf{Projet :} SEPA26 --- Service REST ISO 20022\\
            \textbf{Version :} 1.0.0\\
            \textbf{Date :} Octobre 2026\\
            \textbf{Auteurs :} Umm-Habibah Ouattara et Florian Pépin\\
            \textbf{Technologies :} Spring Boot 4.1.1, Java 25 LTS, Docker, MariaDB
        \end{minipage}%
    }

\end{titlepage}

\newpage

\tableofcontents

\newpage

\section{Introduction}

Ce projet a pour objectif la conception et l'implémentation d'un service web REST complet dédié à la gestion, la validation et la restitution de transactions de prélèvement bancaire SEPA, basées sur la norme internationale \textbf{ISO 20022} (\texttt{pain.008.001.02}).

Le service permet l'ingestion de messages XML, leur validation stricte selon un schéma XML Schema (XSD), leur stockage persistant dans une base relationnelle, leur recherche multicritères, ainsi que leur restitution sous forme de flux XML structurés ou de vues HTML conviviales transformées par feuilles de style XSLT 3.0.

L'application s'accompagne d'une infrastructure moderne et autonome basée sur \textbf{Docker}, \textbf{Docker Compose} et un pipeline d'Intégration Continue (\textbf{GitHub Actions}).

\section{Organisation du Projet et Stack Technique}

Le projet repose sur une séparation nette entre le code applicatif et les descripteurs d'infrastructure, garantissant un environnement modulaire, maintenable et immédiatement opérationnel.

\subsection{Organisation du Dépôt et Arborescence}

Afin d'isoler nettement les outils d'infrastructure du code source applicatif, le dépôt est structuré comme suit :
\begin{itemize}[leftmargin=*]
    \item \textbf{\texttt{app/}} : Application Spring Boot autonome :
    \begin{itemize}
        \item \texttt{src/main/java/} : Code source Java (contrôleurs, services, entités, DTOs).
        \item \texttt{src/main/resources/} : Configuration, templates Thymeleaf, CSS et ressources XML/XSD/XSLT.
        \item \texttt{src/test/} : Suites de tests unitaires et d'intégration.
        \item \texttt{pom.xml} : Descripteur Maven de dépendances et de plugins.
        \item \texttt{checkstyle.xml} : Règles de conformité et d'analyse statique du code (linter).
        \item \texttt{Dockerfile} : Fichier de build multi-stage de l'image Docker.
    \end{itemize}
    \item \textbf{\texttt{doc/}} : Documentation technique et sources LaTeX du présent rapport.
    \item \textbf{\texttt{docker-compose.yml}} : Orchestration locale des conteneurs applicatif et base de données.
    \item \textbf{\texttt{.env} et \texttt{.env.example}} : Paramètres de configuration et variables d'environnement.
    \item \textbf{\texttt{.github/workflows/}} : Pipeline d'Intégration Continue (GitHub Actions).
\end{itemize}

\subsection{Stack Technique}

\begin{itemize}[leftmargin=*]
    \item \textbf{Langage :} Java 25 (LTS).
    \item \textbf{Framework :} Spring Boot 4.1.1 (Spring Web, Spring Data JPA).
    \item \textbf{Gestion des dépendances et Build :} Apache Maven.
    \item \textbf{Qualité de code et Linter :} Maven Checkstyle Plugin 3.6.0.
    \item \textbf{Base de données :} MariaDB 11 en production / Docker, et H2 Database en mémoire pour l'exécution des tests.
    \item \textbf{Traitement XML :} Jakarta XML Binding (JAXB 4.0), Saxon-HE 12.9 (moteur XSLT 3.0), Jackson XML.
    \item \textbf{Moteur de rendu Web :} Thymeleaf.
    \item \textbf{Conteneurisation et CI :} Docker, Docker Compose, GitHub Actions.
\end{itemize}

\newpage

\section{Points d'Accès de l'Application}

L'application expose ses fonctionnalités à travers deux interfaces complémentaires : une API REST ISO 20022 pour les échanges de données XML, et une interface Web Thymeleaf pour la consultation interactive.

\begin{table}[h]
    \centering
    \small
    \renewcommand{\arraystretch}{1.25}
    \begin{tabularx}{\textwidth}{|l|l|l|X|}
        \hline
        \rowcolor{gray!20} \textbf{Méthode} & \textbf{Endpoint} & \textbf{Format} & \textbf{Description} \\
        \hline
        \texttt{GET} & \texttt{/} & HTML & Page d'accueil du service \\
        \hline
        \texttt{GET} & \texttt{/help} & HTML & Documentation interactive des points d'accès \\
        \hline
        \texttt{GET} & \texttt{/transfert} & HTML & Formulaire d'importation et validation de flux XML \\
        \hline
        \texttt{GET} & \texttt{/sepa26/resume/html} & HTML & Tableau de bord des 10 dernières transactions \\
        \hline
        \texttt{GET} & \texttt{/sepa26/resume/xml} & XML & Synthèse XML des 10 derniers documents \\
        \hline
        \texttt{GET} & \texttt{/sepa26/xml/\{id\}} & XML & Fiche complète du document SEPA \\
        \hline
        \texttt{GET} & \texttt{/sepa26/html/\{id\}} & HTML & Rendu stylisé du document via feuille XSLT \\
        \hline
        \texttt{POST} & \texttt{/sepa26/insert} & XML & Insertion d'un flux SEPA (validation XSD) \\
        \hline
        \texttt{DELETE} & \texttt{/sepa26/delete/\{id\}} & XML & Suppression d'une transaction \\
        \hline
        \texttt{GET} & \texttt{/sepa26/search} & XML & Recherche multicritères (date, montant) \\
        \hline
    \end{tabularx}
    \caption{Synthèse des points d'accès de l'application SEPA26}
\end{table}

\section{Usine Logicielle, Conteneurisation et CI/CD}

Pour garantir un déploiement fiable et reproductible sur n'importe quel environnement, le projet s'appuie sur une infrastructure conteneurisée complète.

\subsection{Build Multi-Stage avec Docker}

Le fichier \texttt{Dockerfile} adopte une approche multi-stage optimisant à la fois la sécurité et la légèreté de l'image finale :
\begin{enumerate}[leftmargin=*]
    \item \textbf{Étape de compilation (\textit{Builder}) :} \\
    S'appuie sur l'image \texttt{eclipse-temurin:25-jdk-alpine}. Elle intègre Maven, compile les sources et génère l'archive exécutable \texttt{sepa26server.jar}.
    \item \textbf{Étape d'exécution (\textit{Runtime}) :} \\
    S'appuie sur l'image minimale \texttt{eclipse-temurin:25-jre-alpine}. Seul le binaire compilé y est injecté. L'application est exécutée sous un compte dédié sans privilège (\texttt{spring:spring}) pour une sécurité optimale.
\end{enumerate}

L'image produite est baptisée \texttt{sepa26-app:latest}.

\subsection{Orchestration Locale avec Docker Compose}

Le fichier \texttt{docker-compose.yml} orchestre les deux services nécessaires au fonctionnement complet du projet :
\begin{itemize}[leftmargin=*]
    \item \textbf{Service \texttt{db} :} Conteneur MariaDB 11 avec volume persistant (\texttt{mariadb\_data}) et contrôle de santé (\textit{healthcheck}).
    \item \textbf{Service \texttt{app} :} Conteneur applicatif \texttt{sepa26-app} exposé sur le port \texttt{8100}, démarré automatiquement une fois la base de données saine et prête à recevoir des connexions.
\end{itemize}

Le lancement s'effectue en une commande unique :
\begin{center}
    \texttt{docker compose up --build}
\end{center}

\subsection{Gestion des Variables d'Environnement}

La configuration de l'infrastructure est externalisée via des fichiers d'environnement :
\begin{itemize}[leftmargin=*]
    \item \texttt{.env.example} : Modèle versionné décrivant les variables nécessaires (\texttt{SERVER\_PORT}, \texttt{DB\_URL}, \texttt{DB\_USER}, \texttt{DB\_PASSWORD}, \texttt{DB\_ROOT\_PASSWORD}).
    \item \texttt{.env} : Fichier de configuration local, ignoré par Git pour préserver la confidentialité des paramètres.
\end{itemize}

\subsection{Pipeline d'Intégration Continue (GitHub Actions)}

Un workflow GitHub Actions automatisé (\texttt{.github/workflows/ci.yml}) valide l'intégrité du code à chaque \textit{push} et \textit{pull request} :
\begin{enumerate}[leftmargin=*]
    \item \textbf{Vérification du code (\textit{Checkstyle}) :} Exécution du linter (\texttt{mvn checkstyle:check}) pour valider les règles de nommage et l'absence d'imports superflus.
    \item \textbf{Compilation et tests (\textit{Build et Test}) :} Mise en place de Java 25 (Temurin) avec cache Maven et exécution des tests unitaires (\texttt{mvn clean test}).
    \item \textbf{Validation du conteneur (\textit{Docker Build}) :} Construction de l'image Docker multi-stage à partir du contexte applicatif (\texttt{docker build -t sepa26-app:latest ./app}).
\end{enumerate}

\newpage

\appendix
\section{Annexe : Catalogue Exhaustif des Classes et Composants}

Ce catalogue détaille l'ensemble des classes, interfaces et records développés au sein du projet SEPA26, avec leur localisation, leur type et leur responsabilité métier. Conformément aux principes de conception SOLID et DRY, le modèle de données est entièrement découplé en entités indépendantes et objets de valeur (\texttt{@Embeddable}) réutilisables.

\small
\setlength{\tabcolsep}{6pt}
\renewcommand{\arraystretch}{1.25}

\begin{longtable}{|>{\raggedright\arraybackslash\sffamily\bfseries}p{5.2cm}|>{\raggedright\arraybackslash}p{9.8cm}|}
\hline
\rowcolor{gray!20} \textsf{Composant / Type} & \textsf{Rôle et Responsabilité} \\
\hline
\endfirsthead
\hline
\rowcolor{gray!20} \textsf{Composant / Type} & \textsf{Rôle et Responsabilité} \\
\hline
\endhead

\multicolumn{2}{|l|}{\cellcolor{blue!10}\textbf{Point d'Entrée Principal (\texttt{fr.univrouen.sepa26})}} \\ \hline
Sepa26Application \newline \normalfont\footnotesize\textit{Classe d'application} & Point d'entrée exécutable démarrant le contexte Spring Boot. \\ \hline

\multicolumn{2}{|l|}{\cellcolor{blue!10}\textbf{Contrôleurs REST et Web (\texttt{fr.univrouen.sepa26.controllers})}} \\ \hline
SepaApiController \newline \normalfont\footnotesize\textit{Contrôleur REST} & API REST (@RestController) gérant l'ensemble des flux XML (/sepa26) : consultation, insertion avec validation XSD, suppression et recherche multicritères. \\ \hline
SepaWebController \newline \normalfont\footnotesize\textit{Contrôleur Web} & Contrôleur Web (@Controller) gérant l'interface HTML Thymeleaf : accueil, documentation, transfert, tableau récapitulatif et restitution XSLT. \\ \hline

\multicolumn{2}{|l|}{\cellcolor{blue!10}\textbf{Services Métier et Traitement XML (\texttt{fr.univrouen.sepa26.services})}} \\ \hline
SepaService \newline \normalfont\footnotesize\textit{Service Métier} & Façade métier coordonnant la persistance JPA, le contrôle d'unicité des virements et la délégation aux services XML. \\ \hline
XmlValidation\-Service \newline \normalfont\footnotesize\textit{Interface} & Contrat du service de validation des flux XML selon les schémas XSD. \\ \hline
XmlValidation\-ServiceImpl \newline \normalfont\footnotesize\textit{Implémentation} & Validation XSD avec mise en cache thread-safe du schéma sepa26.xsd. \\ \hline
XmlSerialization\-Service \newline \normalfont\footnotesize\textit{Interface} & Contrat du service de sérialisation et désérialisation XML JAXB. \\ \hline
XmlSerialization\-ServiceImpl \newline \normalfont\footnotesize\textit{Implémentation} & Sérialisation JAXB avec mise en cache du JAXBContext réutilisable. \\ \hline
XsltTransformation\-Service \newline \normalfont\footnotesize\textit{Interface} & Contrat du service de transformation de documents XML en HTML via XSLT. \\ \hline
EndpointDocumentation\-Service \newline \normalfont\footnotesize\textit{Interface} & Contrat fournissant les métadonnées de documentation des points d'accès. \\ \hline
EndpointDocumentation\-ServiceImpl \newline \normalfont\footnotesize\textit{Implémentation} & Documentation descriptive exhaustive des dix points d'accès pour la page d'aide. \\ \hline

\multicolumn{2}{|l|}{\cellcolor{blue!10}\textbf{Utilitaires et Adaptateurs (\texttt{fr.univrouen.sepa26.util})}} \\ \hline
XsltTransformer \newline \normalfont\footnotesize\textit{Composant Saxon} & Moteur Saxon-HE appliquant la feuille de style sepa26.xslt sur les documents. \\ \hline
LocalDateAdapter \newline \normalfont\footnotesize\textit{Adaptateur JAXB} & Conversion bidirectionnelle entre les dates ISO (YYYY-MM-DD) et LocalDate. \\ \hline
LocalDateTime\-Adapter \newline \normalfont\footnotesize\textit{Adaptateur JAXB} & Conversion bidirectionnelle des horodatages ISO et LocalDateTime. \\ \hline

\multicolumn{2}{|l|}{\cellcolor{blue!10}\textbf{Persistance des Données (\texttt{fr.univrouen.sepa26.repository})}} \\ \hline
DocumentRepository \newline \normalfont\footnotesize\textit{Interface Spring Data JPA} & Interface de persistance assurant les opérations CRUD, la recherche par plage de dates et montants, et la vérification par PmtId. \\ \hline

\multicolumn{2}{|l|}{\cellcolor{blue!10}\textbf{Initialisation (\texttt{fr.univrouen.sepa26.init})}} \\ \hline
DataInitializer \newline \normalfont\footnotesize\textit{Composant CommandLineRunner} & Insertion automatique de deux documents types au démarrage si la base est vide. \\ \hline

\multicolumn{2}{|l|}{\cellcolor{blue!10}\textbf{Entités JPA du Domaine (\texttt{fr.univrouen.sepa26.model})}} \\ \hline
Document \newline \normalfont\footnotesize\textit{Entité Racine JPA/JAXB} & Racine d'agrégat (table \texttt{document}) liant l'identifiant technique au bloc d'initiation. \\ \hline
CstmrDrctDbtInitn \newline \normalfont\footnotesize\textit{Entité JPA} & Bloc d'initiation client (table \texttt{cstmr\_direct\_debit\_initiation}), associant l'en-tête \texttt{GrpHdr} et les lots \texttt{PmtInf}. \\ \hline
GrpHdr \newline \normalfont\footnotesize\textit{Entité JPA} & En-tête de message (table \texttt{grp\_hdr}) : identifiant de message, date de création, nombre d'opérations et somme de contrôle. \\ \hline
PmtInf \newline \normalfont\footnotesize\textit{Entité JPA} & Lot d'instructions de paiement (table \texttt{payment\_info}) regroupant le créancier, la date d'échéance et les transactions. \\ \hline
DrctDbtTxInf \newline \normalfont\footnotesize\textit{Entité JPA} & Transaction unitaire élémentaire (table \texttt{transaction\_info}) : montant instruit, débiteur, mandat et compte bancaire. \\ \hline

\multicolumn{2}{|l|}{\cellcolor{blue!10}\textbf{Objets de Valeur Réutilisables (\texttt{fr.univrouen.sepa26.model})}} \\ \hline
Party \newline \normalfont\footnotesize\textit{Composant @Embeddable} & Acteur d'une transaction, mutualisé selon le principe DRY pour l'émetteur (\texttt{InitgPty}), le créancier (\texttt{Cdtr}) et le débiteur (\texttt{Dbtr}). \\ \hline
Account \newline \normalfont\footnotesize\textit{Composant @Embeddable} & Compte bancaire réutilisé pour \texttt{CdtrAcct} et \texttt{DbtrAcct}, encapsulant l'identifiant de compte. \\ \hline
AccountId \newline \normalfont\footnotesize\textit{Composant @Embeddable} & Identifiant de compte bancaire au format normalisé IBAN ou identifiant privé. \\ \hline
Agent \newline \normalfont\footnotesize\textit{Composant @Embeddable} & Établissement financier mandataire réutilisé pour la banque du créancier (\texttt{CdtrAgt}) et du débiteur (\texttt{DbtrAgt}). \\ \hline
FinInstnId \newline \normalfont\footnotesize\textit{Composant @Embeddable} & Identifiant d'institution financière encapsulant le code BIC de l'établissement. \\ \hline
PaymentTypeInfo \newline \normalfont\footnotesize\textit{Composant @Embeddable} & Spécification du type de paiement SEPA (niveau de service, instrument local, séquence). \\ \hline
ServiceLevel \newline \normalfont\footnotesize\textit{Composant @Embeddable} & Niveau de service du protocole bancaire (code \texttt{SEPA}). \\ \hline
LocalInstrument \newline \normalfont\footnotesize\textit{Composant @Embeddable} & Spécification de l'instrument financier local utilisé (code \texttt{SEPA}). \\ \hline
DrctDbtTx \newline \normalfont\footnotesize\textit{Composant @Embeddable} & Informations techniques du prélèvement direct encapsulant le mandat associé. \\ \hline
MndtRltdInf \newline \normalfont\footnotesize\textit{Composant @Embeddable} & Informations du mandat SEPA (Référence Unique de Mandat RUM et date de signature). \\ \hline
InstdAmt \newline \normalfont\footnotesize\textit{Composant @Embeddable} & Montant instruit avec valeur numérique décimale et attribut de devise (\texttt{ccy="EUR"}). \\ \hline
AccountSchemeId \newline \normalfont\footnotesize\textit{Composant @Embeddable} & Enveloppe de l'identifiant de schéma du créancier (ICS). \\ \hline
PrivateId \newline \normalfont\footnotesize\textit{Composant @Embeddable} & Identification privée de l'organisme créancier pour le schéma d'émission. \\ \hline
OtherIdentification \newline \normalfont\footnotesize\textit{Composant @Embeddable} & Identifiant de créancier SEPA (ICS) avec nom de schéma propriétaire. \\ \hline
SchemeName \newline \normalfont\footnotesize\textit{Composant @Embeddable} & Nom du système de schéma propriétaire (\texttt{SEPA}). \\ \hline

\multicolumn{2}{|l|}{\cellcolor{blue!10}\textbf{Objets de Transfert DTO (\texttt{fr.univrouen.sepa26.dto})}} \\ \hline
ValidationResult \newline \normalfont\footnotesize\textit{Record DTO} & Record immutable contenant l'état de validation XSD et l'erreur associée. \\ \hline
ParseResult \newline \normalfont\footnotesize\textit{Record DTO} & Record générique encapsulant le résultat d'une désérialisation XML JAXB. \\ \hline
SepaResponse \newline \normalfont\footnotesize\textit{DTO XML} & Réponse XML standardisée pour les écritures (INSERTED, DELETED, ERROR). \\ \hline
DocumentList \newline \normalfont\footnotesize\textit{DTO XML} & DTO conteneur XML pour la restitution de listes de documents. \\ \hline
SearchResponse \newline \normalfont\footnotesize\textit{DTO XML} & Réponse XML de recherche multicritères (OK, NONE, ERROR). \\ \hline

\end{longtable}

\end{document}
